%% toc: P3 32 32.1,32.2,32.3,32.4
\chapter{Stability and the Forking Calculus}
\label{sec:p03-w03-p03-031-forking-calculus}

Nonforking is the notion of independence that stability theory and its
generalization to simple theories are built on.  A type over a set \(B\)
does not fork over a subset \(A\) when it adds no ``thin'' information
about \(B\) beyond what \(A\) already supplies.  For this to be usable as
an independence relation, it must obey laws: it should be preserved by
automorphisms (invariance), survive shrinking the parameters and enlarging
the base (monotonicity), allow types to be extended to larger parameter
sets (extension), be symmetric, be transitive along chains of bases, and
be determined by a small part of the parameter set (local character).
This chapter states which of these laws hold in every complete theory,
which need simplicity, and which hold in stable theories, and proves each
of them in the annex.  It also proves that every \(\omega\)-stable theory
is superstable, which completes the chain of inclusions between the
stability classes, and collects the stability-spectrum results of the
handbook in one lookup table.

The inputs are a complete first-order theory \(T\) and the definitions of
\cref{sec:p03-stability}: the order property and stability
(\cref{def:p03-stability}), dividing and forking
(\cref{def:p03-dividing-forking}), simplicity in its local-character form
(\cref{def:p03-simple-theory}) and the cardinal stability classes
(\cref{def:p03-stability-spectrum}).  The conclusions are the inclusion
of \cref{thm:p03-w03-p03-031-omega-superstable}, the spectrum entries of
\cref{tab:p03-w03-p03-031-spectrum} and the laws of
\cref{tab:p03-w03-p03-031-calculus}.  The development of the simple case
follows the order of Casanovas's self-contained account
\cite[\S\S~3--5, pp.~266--273]{Casanovas2007Simplicity}, which proves the
simple-theory laws from general model theory together with one
indiscernible-extraction lemma; that lemma is proved here in a version
with a longer input sequence (\cref{thm:p03-w03-p03-031-extraction}).

\section{Stable, superstable, and \texorpdfstring{\(\omega\)}{omega}-stable theories}
\label{sec:p03-w03-p03-031-classes}

The classes of complete theories that this chapter refers to were fixed in
\cref{sec:p03-stability}: stability is the absence of the order property
(\cref{def:p03-stability}); \(\omega\)-stability counts types over
countable sets, and superstability is finite local character of forking
for a stable theory (\cref{def:p03-stability-spectrum}); simplicity is
local character of forking with bases of size at most \(\abs T\)
(\cref{def:p03-simple-theory}).  The inclusions
\[
 \omega\text{-stable}\ \subseteq\ \text{superstable}\ \subseteq\
 \text{stable}\ \subseteq\ \text{simple}
\]
are all proved in the handbook.  The first is
\cref{thm:p03-w03-p03-031-omega-superstable} below, for a language of any
cardinality.  The second holds by \cref{def:p03-stability-spectrum}, which
asks a superstable theory to be stable.  The third is
\cref{thm:p03-w03-p03-035-stable-simple}, proved through a rank that drops
along dividing formulas and is finite when no formula has the order
property.  Each inclusion is strict: \(T_P\) of
\cref{thm:p03-w03-p03-032-tp} is superstable and not \(\omega\)-stable,
\(T_E\) of \cref{thm:p03-w03-p03-032-te} is stable and not superstable,
and the random graph (\cref{thm:p03-w03-p03-031-random-graph}) is simple
and unstable.

The forking laws below are proved under one of three hypotheses: none (any
complete \(T\)), \(T\) simple, or \(T\) stable.  Since
\cref{def:p03-simple-theory} defines simplicity as local character of
nonforking, local character for a stable theory is the third inclusion
above.  The stable forms of all the laws are collected, with that
inclusion as their only additional input, in
\cref{thm:p03-w03-p03-031-stable-calculus}; by the first two inclusions
they apply to superstable and \(\omega\)-stable theories.

\begin{hbconvention}{def:p03-w03-p03-031-monster}{Ambient model, small sets and conjugacy}
Fix a complete theory \(T\) with infinite models in a language of
cardinality \(\abs T\) and a cardinal \(\bar\kappa>\abs T\).  All
elements, tuples and parameter sets are taken from one
\(\bar\kappa\)-saturated model \(\mathfrak C\entails T\)
(every complete type in finitely many variables over a subset of
cardinality below \(\bar\kappa\) that is consistent with \(\mathfrak C\) is
realized in \(\mathfrak C\)); such a model exists under the Axiom of Choice
by \cref{thm:saturated-extension}.  A set or a tuple is \emph{small}
when its cardinality is below \(\bar\kappa\); every parameter set below is
small, and so is every sequence whose length is named.  Tuples are finite
unless their length is stated.

For each application the saturation cardinal is chosen above the parameter
sets and auxiliary sequence lengths used in its proof.  In particular an
application of indiscernible extraction chooses it above the displayed beth
bound.  This is a choice for that application, not an assumption that one
cardinal is closed under all such bounds for all of its smaller cardinals.
One may first pass to a more saturated elementary extension; dividing and
forking for the original parameters are unchanged by
\cref{thm:p03-w03-p03-031-saturation}~(3).

For tuples \(a,b\) of the same length and a set \(A\), write
\(a\equiv_A b\) when \(\operatorname{tp}(a/A)=\operatorname{tp}(b/A)\).
A set of formulas \(\{\varphi(x;b_i):i\in I\}\) is
\emph{\(k\)-inconsistent}, for a positive integer \(k\), when each of its
subsets with \(k\) members is inconsistent with the elementary diagram of
\(\mathfrak C\).  A sequence is \emph{\(A\)-indiscernible} when it is order
indiscernible over \(A\) in the sense of
\cref{def:p03-order-indiscernibles}; lengths of tuples in a sequence may
be infinite but are fixed along the sequence.

Dividing and forking are those of \cref{def:p03-dividing-forking}, read in
\(\mathfrak C\); \cref{thm:p03-w03-p03-031-saturation} shows that their
witnesses can always be found in \(\mathfrak C\).  A \emph{partial type}
\(\pi(x)\) over a set divides or forks over \(A\) when some finite
conjunction of its members does.  A complete type over \(B\supseteq A\)
then forks over \(A\) exactly when it contains a forking formula, as in
\cref{def:p03-dividing-forking}.
\end{hbconvention}

The set \(A\) in ``divides over \(A\)'' is called the \emph{base}.  The
convention fixes where the arguments happen; it adds no hypothesis on
\(T\).  Saturation is used in two ways: to realize types in a small
number of variables over small sets
(\cref{thm:p03-w03-p03-031-saturation}), and to move a configuration by a
type-preserving copy instead of an automorphism.  No homogeneity of
\(\mathfrak C\) is assumed.

\begin{hbdefinition}{def:p03-w03-p03-031-independence}{Independence and Morley sequences}
Let \(A\), \(B\), \(C\) be small sets.  The set \(A\) is \emph{independent
from \(B\) over \(C\)}, written
\[
 A\mathrel{\raisebox{-0.3ex}{\(\smile\)}\mkern-11.5mu\mid\mkern6.5mu}_{C}B,
\]
when for every finite tuple \(a\) from \(A\) the type
\(\operatorname{tp}(a/B\cup C)\) does not fork over \(C\)
\cite[Definition~5.1, p.~271]{Casanovas2007Simplicity}.  A tuple \(a\) is
used in place of the set of its entries.

Let \((X,<)\) be a linear order and \(A\) a small set.  A sequence
\((a_i:i\in X)\) is \emph{\(A\)-independent} when
\(\operatorname{tp}(a_i/A\cup\{a_j:j<i\})\) does not fork over \(A\) for
every \(i\in X\).  A \emph{Morley sequence over \(A\)} is an
\(A\)-independent, \(A\)-indiscernible sequence; it is a Morley sequence
\emph{in} a type \(p\) when each \(a_i\) realizes \(p\)
\cite[Definition~5.2, p.~272]{Casanovas2007Simplicity}.
\end{hbdefinition}

The symbol is a relation between three sets, not a connective of the
object language.  Independence of a sequence is a condition on each member
over its predecessors; it is not assumed that later members are
independent from earlier ones in the opposite direction.  In a simple
theory that reverse condition follows from symmetry
(\cref{thm:p03-w03-p03-031-symmetry}).

\section{Counting types and the stability spectrum}
\label{sec:p03-w03-p03-031-spectrum}

The cardinal classes of \cref{def:p03-stability-spectrum} are defined by
counting types; superstability is defined by finite local character of
forking.  \Cref{sec:p03-w03-p03-032-spectrum} proves the spectrum theorem
and carries stability above \(2^{\abs T}\) to finite local character
through Shelah's degree (\cref{def:p03-w03-p03-032-degree}).  This section
proves that \(\omega\)-stability gives finite degree, hence
superstability, in a language of any cardinality, and then collects the
spectrum results in \cref{tab:p03-w03-p03-031-spectrum}.  The tool is the
Cantor--Bendixson analysis of the space of types over one fixed parameter
set.  Throughout, \(x\) is a tuple of \(m\geq1\) variables and
\(S_m(A)\) is the set of complete types in \(x\) over \(A\).

\begin{hbdefinition}{def:p03-w03-p03-031-cb-rank}{Cantor--Bendixson derivatives of a type space}
Let \(A\) be a small set, \(m\geq1\) and \(X\defeq S_m(A)\).  For a formula
\(\theta(x)\) with parameters from \(A\), let
\([\theta]\defeq\{p\in X:\theta\in p\}\).  Define subsets
\(X^{(\alpha)}\subseteq X\) by induction on the ordinal \(\alpha\):
\begin{enumerate}
  \item \(X^{(0)}\defeq X\);
  \item \(X^{(\alpha+1)}\) is the set of \(p\in X^{(\alpha)}\) such that,
        for every \(\theta\in p\), the set \([\theta]\cap X^{(\alpha)}\)
        contains a type other than \(p\);
  \item \(X^{(\delta)}\defeq\bigcap_{\alpha<\delta}X^{(\alpha)}\) for a
        limit ordinal \(\delta\).
\end{enumerate}
The \emph{Cantor--Bendixson rank} of \(p\in X\) is
\(\operatorname{CB}_A(p)\defeq\alpha\) when
\(p\in X^{(\alpha)}\setminus X^{(\alpha+1)}\), and
\(\operatorname{CB}_A(p)\defeq\infty\) when \(p\in X^{(\alpha)}\) for every
ordinal \(\alpha\).
\index{Cantor--Bendixson rank!of a type space}
\end{hbdefinition}

Passing from \(X^{(\alpha)}\) to \(X^{(\alpha+1)}\) removes the types that
a single one of their formulas singles out among the types that remain:
the isolated points of \(X^{(\alpha)}\) in the topology whose basic open
sets are the sets \([\theta]\).  The derivative involves the one set
\(A\).  Morley's derivative (\cref{def:p03-morley-global-rank}) is a
different object: it takes into account every system into which \(A\) can
be embedded \cite[Definition~2.2, p.~520]{Morley1965}.

\begin{hblemma}{thm:p03-w03-p03-031-cb}{proved=annex:p03-w03-p03-031-cb}{Cantor--Bendixson ranks of types}
Let \(A\) be small, \(m\geq1\) and \(X=S_m(A)\).
\begin{enumerate}
  \item The sets \(X^{(\alpha)}\) decrease as \(\alpha\) increases, and
        there is an ordinal \(\alpha^*\) with
        \(X^{(\beta)}=X^{(\alpha^*)}\) for every \(\beta\geq\alpha^*\).
        Every \(p\in X\) has exactly one rank, and
        \(\operatorname{CB}_A(p)=\infty\) if and only if
        \(p\in X^{(\alpha^*)}\).
  \item If \(\operatorname{CB}_A(p)=\alpha\) is an ordinal, some
        \(\theta\in p\) has \([\theta]\cap X^{(\alpha)}\subseteq\{p\}\), and
        then \([\theta]\cap X^{(\beta)}=\varnothing\) for every
        \(\beta>\alpha\).
  \item If \(X^{(\alpha^*)}\neq\varnothing\), then
        \(\abs X\geq2^{\aleph_0}\).  Hence, if \(X\) is countable, every
        type in \(X\) has an ordinal rank.
  \item Suppose every type in \(X\) has an ordinal rank, and let \(\theta\)
        be a formula in \(x\) with parameters from \(A\) that is realized in
        \(\mathfrak C\).  Then \([\theta]\) contains a type of greatest rank
        \(\alpha_\theta\), and only finitely many types in \([\theta]\) have
        rank \(\alpha_\theta\).
\end{enumerate}
\end{hblemma}

Items~3 and~4 are the two facts used below: countably many types leave no
type of rank \(\infty\), and below one formula the types of top rank are
finite in number.  The countability hypothesis of item~3 is essential.  In
\(T_P\) (\cref{def:p03-w03-p03-032-tp}) there are \(2^{\aleph_0}\) types of
one element over \(\varnothing\) (\cref{thm:p03-w03-p03-032-tp}), and the
lemma then says nothing.

\begin{hbtheorem}{thm:p03-w03-p03-031-omega-superstable}{proved=annex:p03-w03-p03-031-omega-ss}{\(\omega\)-stable theories are superstable}
Let \(T\) be a complete theory, in a language of any cardinality, that is
\(\omega\)-stable in the sense of \cref{def:p03-stability-spectrum}: every
countable parameter set has at most \(\aleph_0\) complete \(1\)-types over
it.  Take the saturation cardinal \(\bar\kappa\) of
\cref{def:p03-w03-p03-031-monster} above \(2^{\abs T}\), so that the
degree of \cref{def:p03-w03-p03-032-degree} is defined.
\begin{enumerate}
  \item \(T\) is stable.
  \item \(\abs{S_m(A)}\leq\aleph_0\) for every countable set \(A\) and every
        \(m\geq1\).
  \item \(\operatorname{Deg}_m(x=x)<\infty\) for every \(m\geq1\).
  \item \(T\) is superstable in the sense of
        \cref{def:p03-stability-spectrum}.  More precisely, if \(p\) is a
        complete type in finitely many variables over a small set \(B\),
        then some finite \(q\subseteq p\) has
        \(\operatorname{Deg}\bigwedge q=\operatorname{Deg}p\), and \(p\)
        does not fork over the finite set of parameters of \(q\).
\end{enumerate}
\end{hbtheorem}

The hypothesis is about countable parameter sets only, and the conclusion
is about types over sets of every size: finite degree is what carries the
information from the countable sets to all of them.  The choice of
\(\bar\kappa\) affects neither hypothesis nor conclusion, since the set of
types over a countable set and nonforking in the sense of
\cref{def:p03-dividing-forking} do not depend on the ambient model
(\cref{thm:p03-w03-p03-031-saturation}~(3)).  The proof finds, from a
formula of infinite degree, a tree of formulas over a countable set whose
siblings are \(k\)-inconsistent (\cref{thm:p03-w03-p03-032-degree-tree}),
and follows a branch along which the greatest Cantor--Bendixson rank
strictly drops at every level; a strictly decreasing sequence of ordinals
cannot be infinite.

When the language is countable there is a second route, through Morley's
analysis.  Under his standing assumption of a complete theory in a
countable language \cite[\S~2, p.~518]{Morley1965}, a theory with
countably many types over every countable set is totally transcendental
and then has at most \(\abs A+\aleph_0\) types over every set \(A\)
\cite[Theorems~2.7 and~2.8, p.~522]{Morley1965}; this is
\cref{thm:p03-morley-countability}.  Such a theory is therefore
\(\lambda\)-stable for every infinite \(\lambda\), in particular for every
\(\lambda\geq2^{\aleph_0}=2^{\abs T}\), and condition~(1) of
\cref{thm:p03-w03-p03-032-superstable} gives superstability.  This route
gives more for countable languages, namely stability in every infinite
cardinal; the proof of \cref{thm:p03-w03-p03-031-omega-superstable} gives
superstability without a countability hypothesis on the language.  The
converse of the theorem fails: \(T_P\) is superstable and not
\(\omega\)-stable (\cref{thm:p03-w03-p03-032-tp}).

\begin{hbexample}{def:p03-w03-p03-031-cb-example}{Ranks of types in an infinite set}
Let \(T_=\) be the theory of an infinite set in the language with equality
only, \(A\) a countably infinite set and \(m=1\).  A permutation of the
universe of \(\mathfrak C\) preserves every formula, so if \(c\) and \(d\)
lie outside a finite set \(\{a_1,\ldots,a_n\}\), the transposition of
\(c\) and \(d\) shows that they satisfy the same formulas with parameters
\(a_1,\ldots,a_n\).  Hence \(S_1(A)\) consists of the types \(p_a\)
containing \(x=a\), for \(a\in A\), and one type \(p_\infty\) of the
elements outside \(A\); there are \(\abs A+1=\aleph_0\) of them, and
\(T_=\) is \(\omega\)-stable.  Each \(p_a\) has rank \(0\):
\([x=a]=\{p_a\}\).  If \(\theta\in p_\infty\) has parameters
\(a_1,\ldots,a_n\), every \(a\in A\setminus\{a_1,\ldots,a_n\}\) satisfies
\(\theta\), so \([\theta]\) contains infinitely many types \(p_a\); thus
\(p_\infty\in X^{(1)}=\{p_\infty\}\), and \(p_\infty\) has rank \(1\) because
\([x=x]\cap X^{(1)}=\{p_\infty\}\).  So \(X^{(2)}=\varnothing\), and
\([x=x]\) has greatest rank \(1\), attained by one type.  For
\(\theta=(x=a)\) the greatest rank is \(0\), again attained once.
\end{hbexample}

\begin{longtable}{@{}p{0.22\textwidth}p{0.33\textwidth}p{0.37\textwidth}@{}}
\caption{The stability spectrum: classes, type counts and where each relation is proved.}\label{tab:p03-w03-p03-031-spectrum}\\
\toprule
Class or condition & Definition or type count & Relations, with proofs and examples \\
\midrule
\endfirsthead
\toprule
Class or condition & Definition or type count & Relations, with proofs and examples \\
\midrule
\endhead
\(\lambda\)-stable, \(\lambda\) infinite & \(\abs{S_1(A)}\leq\lambda\) when
\(\abs A\leq\lambda\) (\cref{def:p03-stability-spectrum}) & then also
\(\abs{S_m(A)}\leq\lambda\) for every \(m\)
(\cref{thm:p03-w03-p03-032-one-to-many}) \\
\addlinespace
stable & no formula has the order property (\cref{def:p03-stability}) &
equivalent to \(\lambda\)-stable for some infinite \(\lambda\), and to
\(\lambda\)-stable for every \(\lambda\) with \(\lambda^{\abs T}=\lambda\);
then \(\abs{S_1(A)}\leq(\abs A+\aleph_0)^{\abs T}\)
(\cref{thm:p03-w03-p03-032-spectrum}) \\
\addlinespace
unstable & some formula has the order property & \(\lambda\)-stable for no
infinite \(\lambda\) (\cref{thm:p03-w03-p03-032-spectrum}); the random
graph (\cref{def:p03-w03-p03-031-random-graph-example}) \\
\addlinespace
stable, \(\lambda^{\abs T}>\lambda\) & not decided by stability &
\(T_E\) is \(\lambda\)-stable exactly when \(\lambda^{\aleph_0}=\lambda\),
so not \(\beth_\omega\)-stable (\cref{def:p03-w03-p03-032-arithmetic},
\cref{thm:p03-w03-p03-032-te}) \\
\addlinespace
\(\lambda\)-stable for all \(\lambda\geq2^{\abs T}\) & Shelah's cardinal
superstability & equivalent to stability in one \(\lambda\) with
\(\lambda^{\aleph_0}>\lambda\geq2^{\aleph_0}\), to stability with
ordinal degree in every finite tuple of variables, and to finite local
character of forking
(\cref{thm:p03-w03-p03-032-superstable}); under these conditions
\(\abs{S_1(A)}\leq\abs A+2^{\abs T}\) \\
\addlinespace
superstable & stable, and every finite-tuple type over \(B\) does not fork
over a finite subset of \(B\) (\cref{def:p03-stability-spectrum}) &
implies stable by definition; \(T_E\) is stable and not superstable
(\cref{thm:p03-w03-p03-032-te}) \\
\addlinespace
\(\omega\)-stable & \(\abs{S_1(A)}\leq\aleph_0\) when \(A\) is countable &
implies stable, finite degree and superstable, any language
(\cref{thm:p03-w03-p03-031-omega-superstable}); for a countable language
\(\abs{S_1(A)}\leq\abs A+\aleph_0\) for every \(A\)
(\cref{thm:p03-morley-countability}); \(T_P\) is superstable and not
\(\omega\)-stable, with \(\abs{S_1(A)}=\abs A+2^{\aleph_0}\)
(\cref{thm:p03-w03-p03-032-tp}) \\
\addlinespace
simple & every finite-tuple type over \(B\) does not fork over some
\(A\subseteq B\) with \(\abs A\leq\abs T\) (\cref{def:p03-simple-theory}) &
stable implies simple (\cref{thm:p03-w03-p03-035-stable-simple}); the random
graph is simple and unstable (\cref{thm:p03-w03-p03-031-random-graph}) \\
\bottomrule
\end{longtable}

In each row the last column says what the class or condition implies,
what implies it, and which theory shows that an implication cannot be
reversed.  The
independence laws that hold in each class are in
\cref{tab:p03-w03-p03-031-calculus}.

\section{Dividing, forking, and nonforking extensions}
\label{sec:p03-w03-p03-031-basic}

\subsection{Indiscernibles and dividing}
\label{sec:p03-w03-p03-031-toolkit}

The following facts hold in every complete theory.  They express dividing
through indiscernible sequences and show how nondividing behaves when the
base is enlarged.

\begin{hblemma}{thm:p03-w03-p03-031-saturation}{proved=annex:p03-w03-p03-031-sat}{Small types and witnesses inside the ambient model}
\begin{enumerate}
  \item Let \(D\) be small and let \(q(x_\alpha:\alpha<\mu)\) be a set of
        formulas with parameters from \(D\) in a small number \(\mu\) of
        variables.  If every finite subset of \(q\) is satisfied in
        \(\mathfrak C\), then \(q\) is realized in \(\mathfrak C\).
  \item If \(c\equiv_A c'\) and \(e\) is any small tuple, there is \(e'\)
        with \(ce\equiv_A c'e'\).
  \item Let \(\mathcal M\preccurlyeq\mathfrak C\) and let \(A\) and the
        parameters of \(\varphi(x;b)\) and \(\psi(x;c)\) lie in
        \(\mathcal M\).  Then \(\varphi(x;b)\) divides over \(A\), and
        \(\psi(x;c)\) forks over \(A\), in the sense of
        \cref{def:p03-dividing-forking} for \(\mathcal M\) if and only if
        this is witnessed by sequences and tuples taken in \(\mathfrak C\).
        The same holds with \(\mathfrak C\) itself in place of
        \(\mathcal M\).
  \item Dividing and forking do not depend on how the parameters are
        listed: if every entry of \(e\) is an entry of the tuple \(e'\)
        and \(\varphi(x;e)\) is regarded as a formula with parameter tuple
        \(e'\) in which the extra entries do not occur, it divides, or
        forks, over \(A\) in one presentation if and only if it does in
        the other.  Likewise adding to \(x\) variables that do not occur
        in the formula changes neither dividing nor forking.
\end{enumerate}
\end{hblemma}

\begin{hblemma}{thm:p03-w03-p03-031-nondividing-indiscernible}{proved=annex:p03-w03-p03-031-ndi}{Nondividing and indiscernible sequences}
Let \(a\) and \(b\) be small tuples and \(A\) a small set.  The following
are equivalent.
\begin{enumerate}
  \item \(\operatorname{tp}(a/A\cup b)\) does not divide over \(A\).
  \item For every infinite \(A\)-indiscernible sequence \(I\) whose first
        member is \(b\), there is \(a'\equiv_{Ab}a\) such that \(I\) is
        \(A\cup a'\)-indiscernible.
  \item For every infinite \(A\)-indiscernible sequence \(I\) whose first
        member is \(b\), there is \(J\equiv_{Ab}I\) that is
        \(A\cup a\)-indiscernible.
\end{enumerate}
\end{hblemma}

This is Casanovas's Lemma~3.1 \cite[p.~267]{Casanovas2007Simplicity},
stated with the indiscernible sequence starting at \(b\), which is the
form \cref{def:p03-dividing-forking} uses.  The lemma converts a
nondividing assumption into a way to move indiscernible sequences while
keeping their type, and it is the engine behind every later result of
this section.

\begin{hblemma}{thm:p03-w03-p03-031-pairs}{proved=annex:p03-w03-p03-031-pairs}{Nondividing of pairs and enlargement of the base}
Let \(A\subseteq B\) be small sets and \(a\), \(b\) small tuples.
\begin{enumerate}
  \item If \(\operatorname{tp}(a/B)\) does not divide over \(A\) and
        \(\operatorname{tp}(b/B\cup a)\) does not divide over \(A\cup a\),
        then \(\operatorname{tp}(ab/B)\) does not divide over \(A\)
        \cite[Proposition~3.4, p.~267]{Casanovas2007Simplicity}.
  \item If \(\varphi(x;a)\) divides over \(A\) and
        \(\operatorname{tp}(b/A\cup a)\) does not divide over \(A\), then
        \(\varphi(x;a)\) divides over \(A\cup b\)
        \cite[Proposition~3.5, p.~267]{Casanovas2007Simplicity}.
\end{enumerate}
\end{hblemma}

\begin{hblemma}{thm:p03-w03-p03-031-extraction}{proved=annex:p03-w03-p03-031-ext}{Indiscernible extraction copying finite types}
Let \(B\) be a small set, \(n_0\) a positive integer,
\(\lambda=\abs T+\abs B\), \(\mu=(2^\lambda)^+\), and let
\(\beth_\alpha(\lambda)\) be defined by \(\beth_0(\lambda)=\lambda\),
\(\beth_{\alpha+1}(\lambda)=2^{\beth_\alpha(\lambda)}\) and suprema at
limits.  For every sequence \((c_i:i<\beth_\mu(\lambda))\) of
\(n_0\)-tuples there is a \(B\)-indiscernible sequence
\((b_i:i<\omega)\) such that, for every \(n<\omega\), some
\(i_0<\cdots<i_n\) satisfy
\[
 b_0\cdots b_n\equiv_B c_{i_0}\cdots c_{i_n}.
\]
\end{hblemma}

Casanovas states this with an input sequence of length \(\mu\) and refers
to Kim's thesis for the proof \cite[Lemma~1.1, p.~264]{Casanovas2007Simplicity}.
The version above asks for a longer input sequence so that it follows
directly from the Erd\H{o}s--Rado theorem (\cref{thm:p06-erdos-rado}); every
later application builds its input sequence by a construction of
arbitrary small length, so the longer requirement costs nothing.  The
conclusion is stronger than the extraction of
\cref{thm:p03-ehrenfeucht-mostowski}: each finite piece of the new
sequence has exactly the type of a piece of the old one, not just the
formulas common to all its increasing pieces.  This is what lets
nonforking pass from the old sequence to the new one.

\subsection{Invariance, monotonicity, extension and existence}
\label{sec:p03-w03-p03-031-laws}

\begin{hbtheorem}{thm:p03-w03-p03-031-basic-calculus}{proved=annex:p03-w03-p03-031-basic}{Invariance, monotonicity and extension in every complete theory}
Let \(T\) be any complete theory and \(A\subseteq B\subseteq C\) small sets.
\begin{enumerate}
  \item \emph{Invariance.}  If \(b\equiv_A b'\), then \(\varphi(x;b)\)
        divides over \(A\) if and only if \(\varphi(x;b')\) does, and
        likewise for forking.  Hence if \(f\) is an elementary map fixing
        \(A\) pointwise and defined on \(B\), a type \(p\) over \(B\) forks
        over \(A\) if and only if its image \(f(p)\) does.
  \item \emph{Dividing implies forking; the forking ideal.}  A dividing
        formula forks.  If \(\psi(x;c)\) implies a finite disjunction of
        formulas each forking over \(A\), then \(\psi(x;c)\) forks over
        \(A\).
  \item \emph{Monotonicity.}  Let \(p\) be a complete type over \(C\) that
        does not fork over \(A\).  Then \(p\mathbin{\upharpoonright}B\) does
        not fork over \(A\); \(p\) does not fork over \(B\); and if
        \(p=\operatorname{tp}(ab/C)\), then \(\operatorname{tp}(a/C)\) does
        not fork over \(A\).  The same three statements hold with dividing
        in place of forking.
  \item \emph{Finite character.}  A complete type over \(C\) forks over
        \(A\) if and only if \(\operatorname{tp}(a/A\cup c)\) forks over
        \(A\) for some realization \(a\) and some finite tuple \(c\) from
        \(C\).
  \item \emph{Extension.}  If a partial type \(\pi(x)\) over \(B\) does
        not fork over \(A\), then some complete type over \(C\) contains
        \(\pi\) and does not fork over \(A\).  In particular a complete
        type over \(B\) that does not fork over \(A\) has an extension to
        \(C\) that does not fork over \(A\).
\end{enumerate}
\end{hbtheorem}

The base conditions matter.  Monotonicity in the base goes in one
direction only: nonforking over a smaller base \(A\) implies nonforking
over a larger base \(B\) inside the domain, not conversely
(\cref{def:p03-w03-p03-031-base-failure}).  Extension in item~5 preserves
nonforking over the same base \(A\) that the given type has; it does not
say that a type over \(A\) does not fork over \(A\).  That last property,
\emph{existence}, is where simplicity first enters.  The extension law is
stated for forking: the forking formulas form an ideal (item~2), and this
is what lets a nonforking type be completed.  For dividing alone no such
closure under disjunction is available in general theories
\cite[p.~267]{Casanovas2007Simplicity}.

\begin{hbtheorem}{thm:p03-w03-p03-031-simple-extension}{proved=annex:p03-w03-p03-031-simp}{Existence and extension in simple theories}
Let \(T\) be simple and \(A\subseteq C\) small sets.
\begin{enumerate}
  \item \emph{Existence.}  No complete type over \(A\) forks over \(A\).
  \item \emph{Extension.}  Every complete type over \(A\) has an
        extension to a complete type over \(C\) that does not fork over
        \(A\).
  \item \emph{Local character.}  For every complete type \(p\) over \(C\)
        in a finite tuple of variables there is \(A_0\subseteq C\) with
        \(\abs{A_0}\leq\abs T\) such that \(p\) does not fork over \(A_0\),
        and then \(p\) does not fork over any \(A'\) with
        \(A_0\subseteq A'\subseteq C\).
\end{enumerate}
\end{hbtheorem}

Item~1 is Casanovas's Corollary~4.1 \cite[p.~271]{Casanovas2007Simplicity};
with the local-character definition of simplicity it is immediate, as
Casanovas observes at the start of his Section~4
\cite[p.~268]{Casanovas2007Simplicity}.  Item~2 is the extension law of
\cref{thm:p03-w03-p03-031-basic-calculus} applied with base \(A\) and
\(B=A\), which existence makes possible.

\begin{hblemma}{thm:p03-w03-p03-031-morley-existence}{proved=annex:p03-w03-p03-031-morley}{Morley sequences}
Let \(T\) be any complete theory.
\begin{enumerate}
  \item If \(A\subseteq B\) are small and \(p\) is a complete type over
        \(B\) in finitely many variables that does not fork over \(A\),
        there is a Morley sequence \((a_i:i<\omega)\) over \(A\) in \(p\)
        that is moreover \(B\)-indiscernible
        \cite[Lemma~5.1, p.~272]{Casanovas2007Simplicity}.
  \item If \((a_i:i\in X)\) is \(A\)-independent and \(Y,Z\subseteq X\)
        satisfy \(y<z\) for all \(y\in Y\), \(z\in Z\), then
        \(\operatorname{tp}((a_i:i\in Z)/A\cup\{a_i:i\in Y\})\) does not
        divide over \(A\) \cite[Remark~5.2, p.~272]{Casanovas2007Simplicity}.
\end{enumerate}
\end{hblemma}

\section{Extension, symmetry, and transitivity}
\label{sec:p03-w03-p03-031-simple}

From here on the theory is simple.  The key step is Kim's lemma: in a
simple theory, one Morley sequence decides whether a formula divides.

\begin{hbtheorem}{thm:p03-w03-p03-031-kim-lemma}{proved=annex:p03-w03-p03-031-kim}{Kim's lemma}
Let \(T\) be simple, \(A\) a small set and \(\varphi(x;a)\) a formula.
Then \(\varphi(x;a)\) divides over \(A\) if and only if for every infinite
Morley sequence \((a_i:i<\omega)\) over \(A\) in
\(\operatorname{tp}(a/A)\) the set \(\{\varphi(x;a_i):i<\omega\}\) is
inconsistent \cite[Proposition~5.3, p.~272]{Casanovas2007Simplicity}.
\end{hbtheorem}

\begin{hbtheorem}{thm:p03-w03-p03-031-forking-dividing}{proved=annex:p03-w03-p03-031-fd}{Forking equals dividing in simple theories}
Let \(T\) be simple.  A partial type \(\pi(x;a)\) forks over a small set
\(A\) if and only if it divides over \(A\)
\cite[Proposition~5.4, p.~272]{Casanovas2007Simplicity}.  In particular a
complete type over \(B\supseteq A\) forks over \(A\) if and only if it
contains a formula that divides over \(A\).
\end{hbtheorem}

\begin{hbtheorem}{thm:p03-w03-p03-031-symmetry}{proved=annex:p03-w03-p03-031-sym}{Symmetry}
Let \(T\) be simple and \(A,B,C\) small.  If
\(A\mathrel{\raisebox{-0.3ex}{\(\smile\)}\mkern-11.5mu\mid\mkern6.5mu}_{C}B\),
then
\(B\mathrel{\raisebox{-0.3ex}{\(\smile\)}\mkern-11.5mu\mid\mkern6.5mu}_{C}A\).
For tuples: \(\operatorname{tp}(a/C\cup b)\) does not fork over \(C\) if
and only if \(\operatorname{tp}(b/C\cup a)\) does not fork over \(C\)
\cite[Proposition~5.5, p.~273]{Casanovas2007Simplicity}.
\end{hbtheorem}

\begin{hbtheorem}{thm:p03-w03-p03-031-transitivity}{proved=annex:p03-w03-p03-031-trans}{Transitivity}
Let \(T\) be simple.
\begin{enumerate}
  \item If \(B\subseteq C\subseteq D\) are small,
        \(A\mathrel{\raisebox{-0.3ex}{\(\smile\)}\mkern-11.5mu\mid\mkern6.5mu}_{B}C\)
        and
        \(A\mathrel{\raisebox{-0.3ex}{\(\smile\)}\mkern-11.5mu\mid\mkern6.5mu}_{C}D\),
        then
        \(A\mathrel{\raisebox{-0.3ex}{\(\smile\)}\mkern-11.5mu\mid\mkern6.5mu}_{B}D\)
        \cite[Proposition~5.6, p.~273]{Casanovas2007Simplicity}.
  \item For small \(A\subseteq B\subseteq C\) and a complete type \(p\)
        over \(C\): \(p\) does not fork over \(A\) if and only if \(p\)
        does not fork over \(B\) and \(p\mathbin{\upharpoonright}B\) does
        not fork over \(A\).
\end{enumerate}
\end{hbtheorem}

In item~2 the direction from left to right holds in every complete theory,
by monotonicity; only the converse uses simplicity.  This is the form of
transitivity used in \cref{thm:p03-nonforking-calculus}.

\begin{longtable}{@{}p{0.2\textwidth}p{0.3\textwidth}p{0.42\textwidth}@{}}
\caption{The nonforking laws and their hypotheses.}\label{tab:p03-w03-p03-031-calculus}\\
\toprule
Law & Hypothesis on \(T\) & Statement and base condition \\
\midrule
\endfirsthead
\toprule
Law & Hypothesis on \(T\) & Statement and base condition \\
\midrule
\endhead
Invariance & none & conjugates over the base \(A\) fork together;
\cref{thm:p03-w03-p03-031-basic-calculus}~(1) \\
Monotonicity & none & shrink parameters, shrink the tuple, or enlarge the
base inside the domain; \cref{thm:p03-w03-p03-031-basic-calculus}~(3) \\
Finite character & none & forking is witnessed by a finite tuple of
parameters; \cref{thm:p03-w03-p03-031-basic-calculus}~(4) \\
Extension & none & a type over \(B\) nonforking over \(A\subseteq B\)
extends to \(C\supseteq B\) nonforking over \(A\);
\cref{thm:p03-w03-p03-031-basic-calculus}~(5) \\
Existence & simple & a type over \(A\) does not fork over \(A\);
\cref{thm:p03-w03-p03-031-simple-extension}~(1) \\
Extension from the domain & simple & every type over \(A\) has an extension
to \(C\supseteq A\) nonforking over \(A\);
\cref{thm:p03-w03-p03-031-simple-extension}~(2) \\
Local character & simple (by definition); stable (proved) & a type over
\(C\) in finitely many variables does not fork over some
\(A_0\subseteq C\) with \(\abs{A_0}\leq\abs T\), finite when \(T\) is
superstable; \cref{thm:p03-w03-p03-031-simple-extension}~(3),
\cref{thm:p03-w03-p03-031-stable-calculus}~(1) \\
Forking \(=\) dividing & simple & over every small base;
\cref{thm:p03-w03-p03-031-forking-dividing} \\
Symmetry & simple & with the same base \(C\) on both sides;
\cref{thm:p03-w03-p03-031-symmetry} \\
Transitivity & simple & along \(B\subseteq C\subseteq D\);
\cref{thm:p03-w03-p03-031-transitivity} \\
\bottomrule
\end{longtable}

The rows with no hypothesis hold in every complete theory, including
unstable and nonsimple ones.  Every row whose hypothesis is ``simple''
holds in particular for stable theories; the next theorem states the
stable forms, with the base condition of each law written out.

\begin{hbtheorem}{thm:p03-w03-p03-031-stable-calculus}{proved=annex:p03-w03-p03-031-stable}{The forking calculus in stable theories}
Let \(T\) be stable, that is, let no formula have the order property
(\cref{def:p03-stability}), and let \(A\subseteq B\subseteq C\) be small
sets.
\begin{enumerate}
  \item \emph{Local character.}  For every complete type \(p\) over \(C\)
        in a finite tuple of variables there is \(A_0\subseteq C\) with
        \(\abs{A_0}\leq\abs T\) such that \(p\) does not fork over \(A_0\);
        then \(p\) does not fork over any \(A'\) with
        \(A_0\subseteq A'\subseteq C\).  If \(T\) is superstable,
        \(A_0\) can be taken finite.
  \item \emph{Invariance.}  If \(f\) is an elementary map fixing \(A\)
        pointwise and defined on \(B\), a complete type \(p\) over \(B\)
        forks over \(A\) if and only if \(f(p)\) forks over \(A\).
  \item \emph{Monotonicity.}  If a complete type \(p\) over \(C\) does not
        fork over \(A\), then \(p\mathbin{\upharpoonright}B\) does not fork
        over \(A\), \(p\) does not fork over \(B\), and if
        \(p=\operatorname{tp}(ab/C)\), then \(\operatorname{tp}(a/C)\) does
        not fork over \(A\).
  \item \emph{Existence and extension.}  No complete type over \(A\) forks
        over \(A\).  Every complete type over \(A\) has an extension to a
        complete type over \(C\) that does not fork over \(A\); more
        generally, a complete type over \(B\) that does not fork over \(A\)
        has an extension to \(C\) that does not fork over \(A\).
  \item \emph{Symmetry.}  For tuples \(a,b\),
        \(\operatorname{tp}(a/C\cup b)\) does not fork over \(C\) if and
        only if \(\operatorname{tp}(b/C\cup a)\) does not fork over \(C\);
        for small sets,
        \(D\mathrel{\raisebox{-0.3ex}{\(\smile\)}\mkern-11.5mu\mid\mkern6.5mu}_{C}E\)
        if and only if
        \(E\mathrel{\raisebox{-0.3ex}{\(\smile\)}\mkern-11.5mu\mid\mkern6.5mu}_{C}D\).
  \item \emph{Transitivity.}  A complete type \(p\) over \(C\) does not
        fork over \(A\) if and only if \(p\) does not fork over \(B\) and
        \(p\mathbin{\upharpoonright}B\) does not fork over \(A\).
  \item \emph{Forking and dividing.}  A formula or partial type forks over
        a small set if and only if it divides over it.
\end{enumerate}
In particular the three laws of \cref{thm:p03-nonforking-calculus} hold,
with the base conditions stated there.
\end{hbtheorem}

Only item~1 uses stability beyond simplicity, and it uses it once: the
proof of \cref{thm:p03-w03-p03-035-stable-simple} chooses, for each
formula \(\delta(x;y)\) and each \(k\), a formula of \(p\) of least local
dividing rank, and the absence of the order property is what makes these
ranks finite.  Items~2 and~3 need no hypothesis at all.  Items~4 to~7 are
the simple-theory laws, which apply because item~1 is the definition of
simplicity.  The bound \(\abs{A_0}\leq\abs T\) in item~1 cannot in general
be improved to a finite bound for stable theories: the theory \(T_E\) of
\cref{thm:p03-w03-p03-032-te} is stable and not superstable, so for it
some type over some set forks over every finite subset of that set.

\begin{hbtheorem}{thm:p03-w03-p03-031-random-graph}{proved=annex:p03-w03-p03-031-rg}{Forking in the random graph}
Let \(T=T_{\rm RG}\) be the theory of the random graph
(\cref{thm:p03-w03-p03-015-graph-types}), \(C\subseteq B\) small sets, and
\(a\) a tuple.  Then \(\operatorname{tp}(a/B)\) forks over \(C\) if and
only if some entry of \(a\) lies in \(B\setminus C\).  Consequently
\(T_{\rm RG}\) satisfies local character with finite bases, so it is
simple.
\end{hbtheorem}

\begin{hbexample}{def:p03-w03-p03-031-random-graph-example}{Reading the calculus in the random graph}
In \(T_{\rm RG}\), take vertices \(a\), \(b\) and a small set \(C\).  By
\cref{thm:p03-w03-p03-031-random-graph} applied with \(B=C\cup\{b\}\),
\(a\) is independent from \(b\) over \(C\) exactly when \(a\neq b\) or
\(a\in C\).  This condition is symmetric in \(a\) and \(b\), as
\cref{thm:p03-w03-p03-031-symmetry} requires.  Adjacency plays no role:
if \(a\) and \(b\) are distinct and adjacent, \(\operatorname{tp}(a/C\cup
b)\) still does not fork over \(C\), although it says something about
\(b\) that \(\operatorname{tp}(a/C)\) does not.  For the type
\(p=\operatorname{tp}(a/B)\) of a single vertex, local character holds with
base \(\varnothing\) when \(a\notin B\) and with base \(\{a\}\) when
\(a\in B\).  Existence is the case \(B=C\): no entry of \(a\) lies in
\(B\setminus B=\varnothing\).  Transitivity along \(C\subseteq C'\subseteq
B\) becomes a statement about sets: if no entry of \(a\) lies in
\(C'\setminus C\) and none lies in \(B\setminus C'\), then none lies in
\(B\setminus C\).

The theory is nevertheless unstable: the edge relation \(E(x;y)\) has the
order property.  Choose vertices in the order
\(a_0,b_0,a_1,b_1,\ldots\), each distinct from all earlier ones; by the
extension axioms \(a_i\) can be taken adjacent to none of the earlier
\(b_j\), and \(b_j\) adjacent to exactly those earlier \(a_i\) with
\(i<j\).  Then \(E(a_i,b_j)\) holds exactly when \(i<j\).  So the
forking calculus of a simple theory need not come from stability.
\end{hbexample}

\begin{hbwarning}{def:p03-w03-p03-031-base-failure}{Monotonicity in the base goes one way}
Nonforking over a smaller base implies nonforking over a larger one
contained in the domain (\cref{thm:p03-w03-p03-031-basic-calculus}~(3)),
but the converse fails even in a simple theory.  In \(T_{\rm RG}\) let
\(a\) be a vertex and \(B=\{a\}\).  By
\cref{thm:p03-w03-p03-031-random-graph}, \(\operatorname{tp}(a/B)\) does
not fork over \(\{a\}\) but forks over \(\varnothing\), because the
formula \(x=a\) divides over \(\varnothing\).  For the same reason the
base in local character cannot in general be taken empty, and in the
symmetry and transitivity laws the base must be the same set on both
sides of the relation being compared.
\end{hbwarning}

\paragraph{Relationships.}
The stable package \cref{thm:p03-nonforking-calculus} consists of
extension from the domain, symmetry for tuples and the two-sided
transitivity law of \cref{thm:p03-w03-p03-031-transitivity}~(2); each is a
simple-theory law in \cref{tab:p03-w03-p03-031-calculus}, and
\cref{thm:p03-w03-p03-031-stable-calculus} proves all three for stable
theories together with local character, invariance and monotonicity.  Superstability
(\cref{def:p03-stability-spectrum}) strengthens local character from
bases of size at most \(\abs T\) to finite bases; the random graph
satisfies the finite form without being stable.  The spectrum results and
the inclusion of the \(\omega\)-stable theories in the superstable ones
are collected in \cref{tab:p03-w03-p03-031-spectrum}.  Morley rank and
\(\omega\)-stability are treated in \cref{sec:p03-morley-rank}.

\paragraph{Simple theories: what is proved here and what is proved
separately.}
The simple-theory laws of this chapter (existence, extension from the
domain, Kim's lemma, forking equals dividing, symmetry and transitivity)
are derived from the local-character definition
\cref{def:p03-simple-theory} alone, with no further hypothesis.  Three
groups of results about simple theories are not proved in this chapter
and are developed in \cref{sec:p03-w03-p03-035-chapter}.
\begin{enumerate}
  \item \emph{Stable theories are simple.}  That a stable theory satisfies
        the local-character definition at all is
        \cref{thm:p03-w03-p03-035-stable-simple}, proved through the local
        dividing rank (\cref{def:p03-w03-p03-035-drank}); local character
        of dividing for stable theories is
        \cref{thm:p03-w03-p03-035-stable-dividing}.
        \Cref{thm:p03-w03-p03-031-stable-calculus}, and with it every
        entry ``stable'' in \cref{tab:p03-w03-p03-031-calculus}, depends on
        \cref{thm:p03-w03-p03-035-stable-simple} and on nothing else from
        that chapter.
  \item \emph{Where the simple theories lie among the other dividing
        lines.}  These results use the independence property and the
        strict order property, which this chapter does not define: the
        strict order property excludes simplicity
        (\cref{thm:p03-w03-p03-035-sop-not-simple}); a simple theory
        without the independence property is stable
        (\cref{thm:p03-w03-p03-035-simple-nip-stable}); the random graph
        is simple, unstable and has the independence property
        (\cref{thm:p03-w03-p03-035-random-graph}); and dense linear orders
        are not simple (\cref{thm:p03-w03-p03-035-dlo}).
  \item \emph{The tree property.}  Shelah's tree property for a formula
        gives the order property
        (\cref{thm:p03-w03-p03-035-tree-order}).  This chapter uses
        simplicity only in the local-character form; the equivalence of
        that form with the definition of simplicity through the absence of
        the tree property is not proved in this chapter, and no result of
        this chapter depends on it.
\end{enumerate}
The refinements that separate the stable case from the simple one inside
the independence calculus are proved in neither chapter: stationarity of
nonforking extensions of types over models, the description of
nonforking through definable types, and the independence theorem over
models, which takes over the role of stationarity in simple theories.  No
result of this chapter uses them.
